\section{Conclusiones del Diseño y Despliegue}

El presente proyecto culmina con el diseño, implementación y validación exitosa
de una red empresarial multi-sede que integra 1.650 usuarios distribuidos en
tres segmentos LAN, con servicios de aplicación centralizados y operación en
doble stack IPv4/IPv6. Se presentan las conclusiones respecto al diseño y
despliegue, así como las lecciones aprendidas y recomendaciones para futuras
implementaciones.

\begin{enumerate}
   \item 
La estrategia de segmentación jerárquica mediante VLSM en IPv4 demostró ser
altamente efectiva para maximizar el aprovechamiento del espacio de
direccionamiento asignado \texttt{172.16.192.0/19}. Al ajustar el prefijo de
cada subred según los requerimientos reales (LAN 1 y 2 con /22, LAN 3 con /24,
enlaces con /30), se minimizó el desperdicio de direcciones mientras se
garantizaba espacio para el crecimiento del 10\% en cada área.


\item
La implementación de capas de seguridad, contraseñas cifradas en dispositivos,
SSHv2 para administración remota, firewall restrictivo en servidores, estableció
una defensa en profundidad. La política de denegar por defecto y permitir solo
puertos críticos redujo significativamente la superficie de ataque, sin afectar
la operación de los servicios requeridos.

\item
Los tres servicios de aplicación (web, correo, chat) se desplegaron exitosamente
en maquinas virtuales con Linux, una con el servicio web y la otra con los 
servicios de correo y chat. La selección de software de código abierto
(Nginx, Postfix, Dovecot) demostró ser confiable y con bajo consumo de recursos,
permitiendo las maquinas con 3 cores y 8GB de RAM manejara la carga de
del despliegue.

La implementación personalizada del servidor de chat en C, desarrollada en etapas
previas del proyecto, se integró sin complicaciones significativas con la
infraestructura de red, demostrando la viabilidad de soluciones personalizadas
en entornos corporativos cuando están bien diseñadas.

\item
Todas las pruebas de validación fueron exitosas en su primer ciclo de ejecución
posterior a la corrección de problemas identificados. Se logró cumplir con todos los criterios de aceptación definidos, incluyendo:

\begin{itemize}
\item Direccionamiento correcto en IPv4 e IPv6
\item Conectividad plena entre todos los segmentos (L3)
\item Disponibilidad de servicios
\item Mecanismos de seguridad operativos
\end{itemize}

\item
El diseño demuestra capacidad de escalabilidad hacia el futuro. El espacio de
direccionamiento IPv4 tiene reserva suficiente para el crecimiento del 10\%
previsto y permite agregar una cuarta LAN antes de agotar el rango
\texttt{172.16.192.0/19}. En IPv6, el prefijo \texttt{2001:dbe:bef9::/48}
ofrece capacidad para decenas de subredes adicionales sin restricción práctica.

La adición de nuevas sedes puede realizarse agregando routers y switches
adicionales sin afectar la configuración existente, dado que OSPF propagará
automáticamente las nuevas rutas.


\end{enumerate}

